\lyseis{7.3}{Πρόσθεση ρητών}
\sol{1} α) $3$ \quad β) $-4$ \quad γ) $-7$ \quad δ) $0$ \quad ε) $0$.
\sol{2} α) $-13$ \quad β) $5$ \quad γ) $-5$ \quad δ) $12$ \quad ε) $-12$ \quad στ) $-1$.
\sol{3} α) $0$ \quad β) $-2{,}5$ \quad γ) $-3{,}2$ \quad δ) $-\frac{3}{4}$ \quad ε) $\frac{1}{6}$ \quad στ) $-\frac{5}{6}+\frac{2}{6}=-\frac{3}{6}=-\frac{1}{2}$.
\sol{4} α) $7$ \quad β) $-3{,}5$ \quad γ) $4$ \quad δ) $\frac{1}{3}$.
\sol{5} α) $21-15=6$ \quad β) $(4+1)+(-1{,}5-2{,}5)=5-4=1$ \quad γ) $10+(-12)=-2$.
\sol{6} α) $0$ \quad β) $\big(37+(-37)\big)+\big((-18)+8\big)=0+(-10)=-10$.
\sol{7} α) Σωστό. β) Λάθος: το πρόσημο είναι του μεγαλύτερου κατ' απόλυτη τιμή. γ) Σωστό. δ) Λάθος: $-7$. ε) Σωστό. στ) Σωστό. ζ) Σωστό ($=3$).
\sol{8} Μεσημέρι $-6+9=3$°. Βράδυ $3+(-5)=-2$°.
\sol{9} Υπόλοιπα: $30,\ 18,\ 26,\ -9,\ 11$. Τελικό υπόλοιπο $11$ €.
\sol{10} $-2+5=3$ και $3+(-4)=-1$. Φτάνει στον όροφο $-1$.
\sol{11} $-18+7=-11$ και $-11-4=-15$. Βρίσκεται στα $-15$ μ.
\sol{12} $8+(-5)+(-7)+12=8$ πόντους.
\sol{13} $-45+80=35$ και $35-20=15$ €.
\sol{14} $12+(-7)+(-9)+4+6=6$.
\sol{15} Όταν είναι αντίθετοι.
\sol{16} Το πρόσημο του προσθετέου με τη μεγαλύτερη απόλυτη τιμή (και $0$ αν είναι αντίθετοι).
\sol{17} Για παράδειγμα $(-2)+(-3)$ και $(-8)+(+3)$.
\sol{18} Όχι: $(-3)+(-2)=-5<-3$.
\sol{19} $(-3)+(-4)=-7$ και $-(3+4)=-7$. Ίσα.
\sol{20} Ελεύθερη απάντηση.
\sol{21} Ελεύθερη απάντηση.
